\chapter{Junk no-go и canonical mapping-cone selector}

Projected-curvature recovery оставил два вопроса: является ли endpoint
projector следствием standard represented differential calculus и почему
relative edge normalization равна \(c=1\). Здесь оба вопроса разделяются.
Нормировка действительно фиксируется внутренней изометрией. Standard junk
quotient, напротив, не выводит нужный projector. Но relative mapping-cone
derivative строит его из единственной канонической координаты трёхузлового
графа.

\section{Сводка результатов gate}

Результат состоит из трёх частей.
\begin{enumerate}
 \item Два edge realizations одной \(\Delta\) имеют одинаковую норму, поэтому
 \(c=1\) не является подгонкой.
 \item Обычный degree-two junk quotient нельзя использовать как объяснение
 endpoint path: в минимальном node calculus он этот path удаляет.
 \item У трёхузловой цепи есть единственная reflection-odd координата
 \((-1,0,1)\). Нормированный коммутатор с ней автоматически убивает
 backtracking curvature и оставляет только endpoint block.
\end{enumerate}
Следовательно, selector теперь выводится внутри relative mapping-cone
action, но ещё не доказано, почему именно этот relative action является
физическим bosonic parent action.

\section{Почему \texorpdfstring{\(c=1\)}{c=1} канонично}

Для
\begin{equation}
 A_\Delta=\Delta,
 \qquad
 B_\Delta=c\Delta^T\varepsilon_2
\end{equation}
операция
\begin{equation}
 T:\Delta\longmapsto\Delta^T\varepsilon_2
\end{equation}
является Frobenius-isometry:
\begin{equation}
 \|T(\Delta)\|_F^2=\|\Delta\|_F^2.
 \label{eq:ps-edge-isometry}
\end{equation}
Transpose сохраняет singular values, а \(\varepsilon_2\) unitary. Кроме
того, endpoint reflection меняет \(4\leftrightarrow\bar4\), а SU(2)
pseudoreality переводит два edge realizations друг в друга. Поэтому единая
reflection/reality-invariant trace--Hodge metric допускает только
\begin{equation}
 |c|=1.
\end{equation}
Оставшаяся фаза удаляется rephasing крайнего module, после чего
\begin{equation}
 \boxed{c=1.}
 \label{eq:ps-canonical-edge-normalization}
\end{equation}
На трёхстах random matrices ошибка identity
\eqref{eq:ps-edge-isometry} меньше \(1.1\times10^{-14}\).

\begin{proposition}[Canonical edge metric]
Если оба edges являются reflection/reality pair одного поля \(\Delta\) и
измеряются одной Hilbert--Schmidt/Hodge metric, relative continuous
normalization отсутствует: \(c=1\) с точностью до removable phase.
\end{proposition}

\section{Контроль standard junk quotient}

Чтобы не приписывать junk calculus желаемое поведение, рассмотрим минимальный
node control
\begin{equation}
 \mathcal A_{\rm node}=\mathbb C^3,
 \qquad
 0\longleftrightarrow1\longleftrightarrow2,
\end{equation}
с любыми ненулевыми weights двух edges. Представленные one-forms и two-forms
вычисляются из
\begin{align}
 \pi_D(e_i\,de_j)&=e_i[D,e_j],\\
 \pi_D(e_i\,de_j\,de_k)&=e_i[D,e_j][D,e_k],
\end{align}
а degree-two junk равен
\begin{equation}
 J_D^2=\pi_D\bigl(d\ker\pi_D|_{\Omega^1}\bigr).
\end{equation}

Exact linear algebra даёт
\begin{equation}
 \rank\pi_D(\Omega^1)=4,
 \qquad
 \rank\pi_D(\Omega^2)=5,
 \qquad
 \rank J_D^2=2.
\end{equation}
Оба endpoint units принадлежат junk:
\begin{equation}
 E_{02},E_{20}\in J_D^2.
 \label{eq:ps-endpoint-in-junk}
\end{equation}
Результат устойчив при complex unequal edge weights; maximum residual в
трёх controls меньше \(1.5\times10^{-15}\).

\begin{nogocriterion}[Junk-projector no-go]
Нельзя утверждать, что standard degree-two junk quotient автоматически
оставляет length-two endpoint curvature. Уже минимальный faithful node
control делает противоположное: endpoint path является junk. Полный
Pati--Salam quotient может иметь дополнительную multiplicity structure, но
до его явного вычисления он не служит derivation найденного projector.
\end{nogocriterion}

\section{Каноническая координата трёхузлового графа}

Incidence Laplacian path graph равен
\begin{equation}
 L_{P_3}=
 \begin{pmatrix}
 1&-1&0\\
 -1&2&-1\\
 0&-1&1
 \end{pmatrix}.
\end{equation}
В reflection-odd sector его наименьший ненулевой eigenvalue равен единице, а
eigenvector единственен с точностью до знака. Нормировка endpoint gap к двум
даёт
\begin{equation}
 h=\operatorname{diag}(-1,0,1).
 \label{eq:ps-canonical-graph-height}
\end{equation}
Это не fitted projector: \(h\) восстанавливается из incidence, reflection и
фиксированной graph-distance normalization.

Определим relative mapping-cone derivative
\begin{equation}
 \delta_h(F)=\frac12[h,F].
 \label{eq:ps-mapping-cone-derivative}
\end{equation}
Для even curvature \(F=\mathcal D_\Delta^2\) diagonal backtracking blocks
коммутируют с \(h\), тогда как blocks \(0\leftrightarrow2\) получают factor
единицу. Поэтому
\begin{equation}
 \|\delta_h(F)\|_F^2=\|F_{02}\|_F^2.
 \label{eq:ps-relative-equals-endpoint}
\end{equation}
При уже выведенном \(c=1\)
\begin{equation}
 \boxed{
 \|\delta_h(\mathcal D_\Delta^2)\|_F^2
 =4\det(\Delta\Delta^\dagger).}
 \label{eq:ps-mapping-cone-kappa-four}
\end{equation}
KO6 doubling и physical half-trace сохраняют coefficient четыре. Все random
tests identities \eqref{eq:ps-relative-equals-endpoint}--
\eqref{eq:ps-mapping-cone-kappa-four} проходят с maximum error меньше
\(10^{-12}\).

\begin{theorem}[Canonical relative selector]
На reflection-symmetric chain \(\bar4\to2_R\to4\) graph incidence однозначно
задаёт normalized relative derivation \(\delta_h\), а reality-invariant
trace--Hodge metric фиксирует \(c=1\). Relative curvature norm без continuous
weights равен exact rank selector
\(4\det(\Delta\Delta^\dagger)\).
\end{theorem}

\section{Точный статус parent action}

Обычная curvature norm остаётся
\begin{equation}
 \|F\|_F^2=\Tr\mathcal D_\Delta^4=4\rho^2
\end{equation}
и не различает rank. Поэтому mapping-cone result не следует из raw spectral
action заменой notation. Он задаёт другой, relative functional
\begin{equation}
 S_{\rm rel}[\Delta]
 =\|\delta_h(\mathcal D_\Delta^2)\|_F^2.
 \label{eq:ps-relative-parent-action}
\end{equation}
Следующий gate должен вывести выбор \(S_{\rm rel}\), а не полного
\(\|F\|^2\), из boundary/mapping-cone principle или заранее объявленного
relative observable. После этого требуется решить physical versus auxiliary
status двух endpoint modules и пересчитать mixed Hessian и gauge-running
ledger.

\begin{theorem}[Уточнённое переоткрытие]
Происхождение projector и normalization теперь локализовано.
Standard junk route закрыт, \(c=1\) выведено, а endpoint projector
является commutator с канонической graph coordinate. Незакрытым
остаётся один концептуальный выбор: почему физический parent action
измеряет relative mapping-cone curvature, а не full curvature.
\end{theorem}

\section{Литература и воспроизводимость}

Определение degree-two junk сверено с T.~Sch\"ucker,
arXiv:hep-th/9312186. Superconnection curvature как even operator
\(F=\mathbb A^2\) сверена с G.~Roepstorff, arXiv:hep-th/9801040.
Graph Laplacian/Dirac и distance-function framework сверены с M.~Requardt,
arXiv:math-ph/0001026. Relative mapping-cone specialization
\eqref{eq:ps-mapping-cone-derivative} является конструкцией проекта, а не
утверждением, найденным в этих источниках.

Machine audit сохранён в
\path{s2t_v4_pati_salam_junk_mapping_cone_gate.py} и
\path{s2t_v4_pati_salam_junk_mapping_cone_gate_results.json}.